我衷心希望软件开发者能从本书中获益。开发一个应用涉及收集和分析需求、运用良好的设计原则，并在合适的地方依照设计模式来为代码建模。本书的组织结构正是遵循这一递进过程。

\subsection{本书适合谁阅读？}

如果你想成为一名顶尖程序员，能够开发出满足需求、更加可靠、灵活且可维护的设计良好的应用，那么本书就是为你而写。设计良好的应用 bug 更少，生产成本更低，运行时性能更好。我们将迎战良好软件设计的两大敌人：变化与复杂度。你会为自己设计和开发的应用感到自豪。只要你至少是一名初级的 C++ 程序员并了解面向对象编程的基础知识，就能从本书中收获最多。

\subsection{本书的组织结构：路线图}

本书共 16 章，分为五个部分。第 1 部分“引言\emph{,}”介绍软件设计和一种开发方法：

\begin{itemize}
\item 第 1 章讨论什么是软件设计，并包含若干设计示例。
\item 第 2 章展示如何通过迭代式开发实现良好的设计。
\end{itemize}

第 2 部分“设计出正确的应用”讨论需求与良好的类设计：

\begin{itemize}
\item 第 3 章讲解如何分析需求，并介绍 UML 图。
\item 第 4 章讲解良好的类设计。
\end{itemize}

第 3 部分“把应用设计对”包含良好设计原则的示例，如封装、松耦合、面向接口编程以及隐藏实现：

\begin{itemize}
\item 第 5 章解释为什么隐藏类的实现很重要。
\item 第 6 章解释为什么我们不应该编写让用户感到意外的代码。
\item 第 7 章讲解如何正确设计子类。
\end{itemize}

第 4 部分“设计模式解决应用架构问题”展示如何应用经过业界验证的设计模式：

\begin{itemize}
\item 第 8 章讲解模板方法和策略设计模式。
\item 第 9 章讲解工厂方法和抽象工厂设计模式。
\item 第 10 章讲解适配器和外观设计模式。
\item 第 11 章讲解迭代器和访问者设计模式。
\item 第 12 章讲解观察者设计模式。
\item 第 13 章讲解状态设计模式。
\item 第 14 章讲解单例、组合和装饰器设计模式。
\end{itemize}

第 5 部分“其他设计技术”探讨递归、回溯和多线程：

\begin{itemize}
\item 第 15 章讲解用递归和回溯设计解决方案。
\item 第 16 章讲解多线程程序。
\end{itemize}

我把本书设计成按顺序阅读，尤其是第 1 部分到第 4 部分。具体来说，第 4 部分的各章会引用第 2 部分和第 3 部分的各章所涵盖的内容。良好的软件设计并不仅仅是在这里或那里使用一些规定的编码技巧，而在于各种设计原则和模式相互依存、协同作用。

少数章节被标记为可选。我收录它们是为了完整，也因为它们的主题很有意思。不过，理解本书其余内容并不需要阅读这些章节。

\subsection{关于代码}

学生（以及读者）通过大量示例学习的效果最好。因此，我收录了许多设计不良的程序的示例，并说明了它们为什么不好，以及如何把它们改造成设计良好的程序。

本书的示例程序使用 2017 版 C++，不过其中几乎所有程序都能用更早版本的语言编译和运行。除第 3 章外，每一章都有示例程序。这些程序只使用有限的 C++ 特性，以便设计概念清晰明了，并让你能够把这些概念应用到其他面向对象语言中。

C++ 是一门非常流行但也很复杂的编程语言。它源于 C 语言，而 C 语言通常用于底层系统编程，但 C++ 也用于高层应用开发。这门语言保留了许多略显繁琐的特性，这些特性旨在帮助编译器生成更优化的代码，并让程序员循规蹈矩。我并不想过度使用这类特性。一个例子是关键字 \texttt{const}，它的用法被严重重载，可能相当令人困惑。在我的示例中，我力求以直接、一致的方式使用 \texttt{const} 和其他语言特性。

本书既包含编号代码清单中的源代码，也包含与正文排在同一行内的源代码。两种情况下，源代码都采用\texttt{像这样的等宽字体}排印，以与普通文本区分开。有时，代码还会以\texttt{\textbf{粗}\textbf{体}}显示，以突出显示相对本章前面步骤发生了变化的代码，例如新功能追加到某行既有代码时。

在许多情况下，原始源代码经过了重新排版：添加换行并调整缩进，以适应书中可用的页面空间。在少数情况下，即使这样也不够，代码清单会包含续行标记（↪）。此外，当正文中对代码进行说明时，代码清单中的注释往往已被删除。许多代码清单都配有代码标注，以突出显示重要概念。

你可以从本书的 liveBook（在线）版本获取可执行的代码片段：\href{https://livebook.manning.com/book/object-oriented-software-design-in-c-plus-plus}{\nolinkurl{https://livebook.manning.com/book/object-oriented-software-design-in-c-plus-plus}}。本书示例的完整代码可从 Manning 网站下载：\href{https://www.manning.com/books/object-oriented-software-design-in-c-plus-plus}{\nolinkurl{https://www.manning.com/books/object-oriented-software-design-in-c-plus-plus}}，也可从 GitHub 下载：\href{https://github.com/RonMakBooks/SoftwareDesignCpp}{\nolinkurl{https://github.com/RonMakBooks/SoftwareDesignCpp}}。在那里，你能找到在 Linux、macOS 和 Microsoft Windows 平台上编译和运行这些程序的说明。作为额外福利，你还能找到用于基于本书各章进行授课的幻灯片。

\subsection{liveBook 讨论论坛}

购买\emph{《面向对象软件设计：C++ 实践》}（Object-Oriented Software Design in C++）即可免费访问 liveBook——Manning 的在线阅读平台。借助 liveBook 独有的讨论功能，你可以把评论附加到整本书上，也可以附加到特定的章节或段落上。为自己做笔记、提出并回答技术问题，以及从作者和其他用户那里获得帮助，都轻而易举。要访问论坛，请前往 \href{https://livebook.manning.com/book/object-oriented-software-design-in-c-plus-plus/discussion}{\nolinkurl{https://livebook.manning.com/book/object-oriented-software-design-in-c-plus-plus/discussion}}。你也可以在 \href{https://livebook.manning.com/discussion}{\nolinkurl{https://livebook.manning.com/discussion}} 了解有关 Manning 论坛及其行为准则的更多信息。

Manning 对读者的承诺是提供一个场所，让读者之间以及读者与作者之间能够进行有意义的对话。这并不是承诺作者会参与多少，作者对论坛的贡献仍然是自愿的（而且是无偿的）。我们建议你试着向作者提一些有挑战性的问题，免得他们的兴趣转移了！只要本书仍在印行，论坛以及以往讨论的存档就可以从出版商的网站上访问。
